% !TeX root = ../main.tex
\section{Mô tả mô hình (Model Descriptions)}
Chọn một VLM nhỏ và bốn OCR chuyên dụng để khảo sát khả năng sinh văn bản so với
pipeline phát hiện--nhận dạng. M1--M4 chạy trên một Tesla T4 của Kaggle;
M5 chạy CPU Xeon 2,00 GHz (giới hạn OpenMP 2 luồng).

\begingroup\fontsize{9.2}{10.3}\selectfont
\setlength{\tabcolsep}{3pt}\renewcommand{\arraystretch}{1.05}
\begin{tabularx}{\linewidth}{@{}l>{\hsize=0.95\hsize\linewidth=\hsize}Yl>{\hsize=1.05\hsize\linewidth=\hsize}Y@{}}
\toprule
\textbf{Mã} & \textbf{Mô hình pretrained} & \textbf{Quy mô$^*$} & \textbf{Cấu hình chạy} \\
\midrule
M1 & Florence-2-base~\cite{florence} & $\sim$230M tham số & FP16; $768^2$; \texttt{<OCR>}; 1.024 token; 3 beams; no-repeat 3-gram \\
M2 & EasyOCR 1.7.2: CRAFT + english\_g2~\cite{easyocr} & 24,55M; 93,74 MiB & FP32; canvas 2.560; greedy; batch crop 16 \\
M3 & docTR 0.11.0: DB/CRNN MobileNetV3~\cite{doctr} & 6,26M; 24,24 MiB & FP32; detector $1024^2$; crop $32\!\times\!128$; batch 128 \\
M4 & PaddleOCR 3.4.1: PP-OCRv5 mobile det + en rec~\cite{paddle} & 11,89 MiB & FP32; cạnh dài tối đa 960; batch 16; tắt xoay, unwarp, TensorRT \\
M5 & Tesseract 5.5.4: eng tessdata\_fast~\cite{tesseract} & 3,92 MiB & LSTM số nguyên; OEM 1, PSM 3; không resize; tắt học thích nghi \\
\bottomrule
\end{tabularx}
\endgroup
{\footnotesize $^*$M1 là số tham số danh định; M2/M3 đếm cả hai mạng (docTR dùng
DB MobileNetV3 Large, CRNN MobileNetV3 Small). MiB là tổng tệp trọng số đã lưu;
M5 gồm cả từ điển, không suy ra số tham số từ dung lượng. M4 dùng Paddle 3.2.0,
PaddleX 3.4.0; checkpoint/revision, hash và phiên bản được lưu trong cấu hình.\par}
